\section{Заключение}
\label{sec:zaklyuchenie}

Първо, квалифицираният електронен подпис е единственият електронен подпис, който е
равностоен на саморъчния по силата на закона (чл.~25, пар.~2 eIDAS). Обикновеният и
усъвършенстваният подпис имат това действие само ако страните са го уговорили (чл.~13,
ал.~4 ЗЕДЕУУ). От тази разлика зависят както формалната доказателствена сила на документа
по чл.~180 ГПК, така и разпределението на доказателствената тежест при оспорването му.

Второ, валидирането доказва, че документът не е променян и че е подписан с ключа, свързан
с удостоверението, но не доказва кой човек е дал командата за подписване. Правната
равностойност на КЕП стъпва на предположението, че ключът е единствено под контрола на
титуляря. При дистанционното подписване ключът се пази от доставчика и това предположение
зависи от сигурността на неговите системи.

Трето, когато документът носи валиден КЕП на оспорващия, тежестта да докаже неистинността
му е негова (чл.~193, ал.~3 ГПК). Преди това обаче страната, която представя документа,
трябва да установи, че подписът е валиден КЕП, а при спор съдът е длъжен да провери всички
условия за това. Чл.~18а ЗЕДЕУУ изключва оспорването спрямо адресата, ако титулярят сам е
дал достъп до подписа си или документът е изпратен от автоматична система.

Четвърто, доказателствата за неистинността обикновено са у доставчика. Той е длъжен да ги
пази за целите на съдебните производства (чл.~24, пар.~2, б.~„з“ eIDAS), а съдът може да
го задължи да ги представи (чл.~192 ГПК). Затова тежестта на оспорващия, макар и трудна,
може да бъде изпълнена.

Пето, рискът от злоупотреби е разпределен между титуляря, адресата и доставчика. Казусът
с „Автогара-Хасково“~ЕООД показва, че задължението на доставчика да провери самоличността
и представителната власт се прилага стриктно. Презумпцията за вина на квалифицирания
доставчик (чл.~13, пар.~1 eIDAS) е основната защита на лицето, чиято самоличност е
използвана за издаване на удостоверение.

Работата позволява и три предложения \textit{de lege ferenda}. ЗЕДЕУУ не е изменян от
2019~г. и трябва да бъде съобразен с Регламент (ЕС) 2024/1183, включително с уредбата на
дистанционното подписване. В ГПК би било полезно изрично правило, че при оспорване на
документ с КЕП съдът може и служебно да изиска от доставчика данните за проверката на
самоличността и за командите за подписване. Накрая, чл.~18а, ал.~1, т.~2 ЗЕДЕУУ заслужава
уточнение дали предоставянето на данните за достъп до мобилно приложение за дистанционно
подписване е „предоставяне на достъп до начина на идентифициране“.
